% SOURCE_AUTHORITY: sga5_fr_workpass.tex, lines 15276--15286 (LF-normalized unit).
% SOURCE_SHA256: 791F4EFFC5E02832D5D77ED03518C8156D6F07E4C8238B03545DB93D883FBB28
Demonstração do lema 1: pode-se (substituindo, se necessário, $L$ por um complexo quase-isomorfo) supor que o complexo $L$, e portanto também $T(L)$, seja projetivo em cada grau, o que suporemos no que segue; pode-se, por conseguinte, supor que o isomorfismo $u$ de $D^-(A_0)$ provenha de um quase-isomorfismo $T(L)\to M$, denotado igualmente por $u$. Primeiro nos reduziremos ao caso em que $u$ seja sobrejetivo. O complexo $N_0$ definido por
\[
N_0^i=M^i\oplus M^{i-1},
\quad
 d^i(x,y)=(d^i x,x-d^{i-1}y)
\]
($x\in M^i$, $y\in M^{i-1}$) é manifestamente acíclico e projetivo em cada grau; define-se um epimorfismo de complexos $N_0\to M$ mediante a projeção $N_0^i\to M^i$ em grau $i$. O corolário do lema 3 permite levantar $N_0$ a um complexo $N$ acíclico, limitado à direita e projetivo em cada grau. Pode-se então substituir $L$ por $L\oplus N$, que lhe é quase-isomorfo (pois $N$ é acíclico) e projetivo em cada grau; ao mesmo tempo, substitui-se $u$ pelo quase-isomorfismo
\[
T(L\oplus N)=T(L)\oplus N_0\longrightarrow M
\]
definido por $u$ sobre $T(L)$ e pelo epimorfismo $N_0\to M$ sobre $N_0$; é claro que este último quase-isomorfismo é sobrejetivo.
